% =====================================================================
%  ANEXO G. Matriz de trazabilidad
% =====================================================================
\section{Matriz de trazabilidad}
\label{anx:trazabilidad}

\begin{notanormativa}[Para qué sirve la trazabilidad]
La cadena debe poder recorrerse en ambos sentidos: desde un requisito hasta el resultado que lo
evidencia, y desde un incidente hasta el requisito y el riesgo que justificaron probarlo.
\end{notanormativa}

\begin{observacion}[Trazabilidad sin cadena definida]
\obsproblema{La plantilla anterior pedía «elaborar una tabla de trazabilidad» para los elementos
de prueba, sin indicar qué debía encadenar ni con qué identificadores.}
\obscambio{Matriz de nueve columnas —base, riesgo, condición, elemento de cobertura, caso,
procedimiento, ejecución, evidencia e incidente—, con la Figura~\ref{fig:trazabilidad} que
señala dónde se registra cada eslabón y un resumen de cobertura de la base de prueba.}
\obsfuente{OBS §6}
\end{observacion}

% =====================================================================
%  diagrams/cadena-trazabilidad.tex
%  Figura: cadena de trazabilidad que la plantilla permite representar,
%          con el apartado o anexo donde se registra cada eslabón.
% =====================================================================
\begin{figure}[htbp]
\centering
\begin{tikzpicture}[
  font=\scriptsize,
  esl/.style ={draw=uvazul, thick, fill=uvazul!8, rounded corners=2pt,
               text width=2.45cm, align=center, minimum height=0.95cm, inner sep=3pt},
  ori/.style ={font=\tiny\itshape, text=gray!65, align=center, text width=2.6cm},
  fl/.style  ={-{Latex[length=4pt]}, draw=uvazul!85, thick}
]

\node[esl] (r1) at (-5.6,0)  {\textbf{Base de prueba}\\ BP-xx};
\node[esl] (r2) at (-2.8,0)  {\textbf{Riesgo}\\ RP-xx};
\node[esl] (r3) at (0,0)     {\textbf{Condición de prueba}\\ COND-xx};
\node[esl] (r4) at (2.8,0)   {\textbf{Elemento de cobertura}};
\node[esl] (r5) at (5.6,0)   {\textbf{Caso de prueba}\\ CP-xxx};

% La segunda fila se dispone de derecha a izquierda («en boustrophedon»)
% para que el enlace entre CP-xxx y PP-xxx sea un salto vertical corto y
% no una línea que atraviese todo el diagrama.
\node[esl] (r6) at (5.6,-2.6)  {\textbf{Procedimiento}\\ PP-xxx};
\node[esl] (r7) at (2.8,-2.6)  {\textbf{Ejecución}\\ EJ-xxx};
\node[esl] (r8) at (0,-2.6)    {\textbf{Evidencia}\\ EVID-xxx};
\node[esl] (r9) at (-2.8,-2.6) {\textbf{Incidente}\\ INC-xx};

\foreach \a/\b in {r1/r2, r2/r3, r3/r4, r4/r5, r6/r7, r7/r8, r8/r9} { \draw[fl] (\a) -- (\b); }
\draw[fl] (r5.south) -- (r6.north);

\node[ori] at (-5.6,-1.0) {Cuadro~\ref{tab:base-prueba}};
\node[ori] at (-2.8,-1.0) {Cuadro~\ref{tab:riesgos-producto}};
\node[ori] at (0,-1.0)    {Cuadro~\ref{tab:condiciones}};
\node[ori] at (2.8,-1.0)  {Cuadro~\ref{tab:condiciones}};
\node[ori] at (5.6,-1.0)  {Anexo A};

\node[ori] at (5.6,-3.6)  {Anexo B};
\node[ori] at (2.8,-3.6)  {Anexo C};
\node[ori] at (0,-3.6)    {Anexo C};
\node[ori] at (-2.8,-3.6) {Anexo D};

\node[align=left, font=\scriptsize, text width=15.2cm] at (0,-4.8)
  {\textbf{Cómo leer la figura.} Cada eslabón se registra en el cuadro o anexo indicado bajo él.
   La cadena se recorre en ambos sentidos: hacia adelante, para comprobar que un requisito quedó
   verificado; hacia atrás, para saber qué requisito y qué riesgo justificaron una prueba que
   reveló un incidente. No todos los proyectos necesitan todos los eslabones; la plantilla
   permite representarlos, no obliga a usarlos.};

\end{tikzpicture}
\caption[Cadena de trazabilidad representable en la plantilla]%
{Cadena de trazabilidad desde la base de prueba hasta el incidente, con el cuadro o anexo en
que se registra cada eslabón.}
\label{fig:trazabilidad}
\end{figure}


La Figura~\ref{fig:trazabilidad} muestra la cadena que esta plantilla permite representar y el
apartado o anexo donde se registra cada eslabón. No todos los proyectos necesitan todos los
eslabones: un proyecto sin análisis de riesgo formal, o sin elementos de cobertura explícitos,
puede dejar esas columnas vacías. Lo que la plantilla garantiza es que sean \emph{expresables}.

\begin{instruccion}
Complete una fila por cada relación que quiera poder demostrar. Los identificadores provienen de
los cuadros ya definidos: base de prueba (Cuadro~\ref{tab:base-prueba}), riesgos
(Cuadro~\ref{tab:riesgos-producto}), condiciones (Cuadro~\ref{tab:condiciones}), casos
(Cuadro~\ref{tab:indice-casos}), procedimientos (Cuadro~\ref{tab:indice-procedimientos}),
ejecuciones (Cuadro~\ref{tab:registro-ejecuciones}) e incidentes
(Cuadro~\ref{tab:registro-incidentes}).

Use la matriz también como herramienta de control: una fila de la base de prueba sin caso
asociado indica una brecha de cobertura; un caso sin base ni riesgo indica esfuerzo que nadie
pidió.
\end{instruccion}

\begin{footnotesize}
\begin{longtable}{|C{1.35cm}|C{1.35cm}|C{1.45cm}|C{1.8cm}|C{1.3cm}|C{1.3cm}|C{1.3cm}|C{1.45cm}|C{1.3cm}|}
\caption{Matriz de trazabilidad de la base de prueba al incidente.}
\label{tab:trazabilidad} \\
\hline
\filaenc \textbf{Base} & \textbf{Riesgo} & \textbf{Condición} & \textbf{Elemento de cobertura} & \textbf{Caso} & \textbf{Proc.} & \textbf{Ejec.} & \textbf{Evidencia} & \textbf{Incid.} \\ \hline
\endfirsthead
\hline
\filaenc \textbf{Base} & \textbf{Riesgo} & \textbf{Condición} & \textbf{Elemento de cobertura} & \textbf{Caso} & \textbf{Proc.} & \textbf{Ejec.} & \textbf{Evidencia} & \textbf{Incid.} \\ \hline
\endhead
\hline \multicolumn{9}{|r|}{\textit{continúa en la página siguiente}} \\ \hline
\endfoot
\hline
\endlastfoot
BP-01 & RP-01 & COND-01 & Partición P3 & CP-007 & PP-001 & EJ-014 & EVID-014 & INC-05 \\ \hline
\campo{} & \campo{} & \campo{} & \campo{} & \campo{} & \campo{} & \campo{} & \campo{} & \campo{} \\ \hline
 & & & & & & & & \\ \hline
 & & & & & & & & \\ \hline
 & & & & & & & & \\ \hline
 & & & & & & & & \\ \hline
\end{longtable}
\end{footnotesize}

La primera fila del Cuadro~\ref{tab:trazabilidad} está rellena a modo de ejemplo y reproduce la
cadena completa: el requisito BP-01, motivado por el riesgo RP-01, generó la
condición COND-01, cuya partición P3 se ejercita mediante el caso CP-007, incluido en el
procedimiento PP-001; su ejecución EJ-014 produjo la evidencia EVID-014 y dio origen al
incidente INC-05. Sustituya esa fila por los datos reales del proyecto.

\begin{instruccion}
Si el volumen de filas hace inmanejable esta matriz en el documento, manténgala en una hoja de
cálculo o herramienta y deje aquí la referencia a su ubicación y responsable, igual que con el
resto de los registros. Lo que no es admisible es no poder reconstruir la cadena.
\end{instruccion}

\subsection*{Cobertura de la base de prueba}

\begin{instruccion}
Este cuadro resume la matriz anterior desde el punto de vista de la base de prueba. Es la
fuente del dato de la medida COB-01 del Cuadro~\ref{tab:cobertura} y sirve para detectar
elementos de la base que quedaron sin ejercitar.

Las dos últimas columnas responden a preguntas distintas y \textbf{no deben fundirse}:
«¿Cubierto?» dice si el elemento llegó a ejercitarse; «Veredictos» dice cómo se comportó. Un
requisito ejercitado por casos que fallaron está \emph{cubierto} y además tiene un
\emph{resultado adverso}: cuenta en el numerador de COB-01 y, a la vez, en los incidentes del
\ref{anx:incidentes}.
\end{instruccion}

\begin{table}[H]
\centering
\caption{Resumen de cobertura de la base de prueba y de sus resultados.}
\label{tab:cobertura-base}
\begin{tabularx}{\textwidth}{|C{1.5cm}|Y|C{1.9cm}|C{1.9cm}|C{2.2cm}|C{2.3cm}|}
  \hline
  \filaenc \textbf{Base} & \textbf{Descripción} & \textbf{Casos asociados} & \textbf{Casos ejecutados} & \textbf{¿Cubierto? (COB-01)} & \textbf{Veredictos (apr. / no apr.)} \\ \hline
  \campo{BP-01} & \campo{Descripción} & \campo{3} & \campo{3} & \campo{Sí} & \campo{3 / 0} \\ \hline
  \campo{BP-02} & \campo{Descripción} & \campo{2} & \campo{2} & \campo{Sí} & \campo{0 / 2} \\ \hline
  \campo{BP-03} & \campo{Descripción} & \campo{0} & \campo{0} & \campo{No — brecha} & \campo{---} \\ \hline
  & & & & & \\ \hline
\end{tabularx}
\notatabla{«¿Cubierto?» es \textit{Sí} en cuanto el elemento tiene al menos un caso ejecutado, cualquiera que sea el veredicto. Sólo las columnas de casos ejecutados y de veredictos cambian al repetir una prueba; la de cobertura no retrocede.}
\end{table}

El Cuadro~\ref{tab:cobertura-base} hace visibles dos situaciones que conviene no confundir. La
fila BP-03 es una \textbf{brecha de cobertura}: un elemento de la base que quedó sin ejercitar
y que, si no se justifica como exclusión en el Cuadro~\ref{tab:exclusiones}, debe aparecer
entre los riesgos residuales del apartado~\ref{sec:residuales}. La fila BP-02, en cambio,
está plenamente cubierta y aun así arroja un \textbf{resultado adverso}: sus dos casos
fallaron. Fundir ambas columnas en una sola haría desaparecer una de las dos situaciones, y
contar BP-02 como «no cubierta» ocultaría que la prueba sí se realizó y encontró un defecto.
